% !TeX root = ../main.tex
\section{Cơ sở hình thành ý tưởng và lý do sản phẩm có tính khả thi}

Theo thông tin trong bài đăng được cung cấp, Chính phủ đã xác định \textbf{game là một trong 6 ngành công nghiệp văn hóa trọng điểm}. Bên cạnh đó, việc một giải đấu thể thao điện tử cấp đội tuyển quốc gia quy tụ nhiều quốc gia Đông Nam Á được tổ chức tại Việt Nam cho thấy eSports không còn chỉ là hoạt động giải trí cá nhân, mà đang dần được nhìn nhận như một lĩnh vực có tiềm năng phát triển chính thức.

Điều này tạo ra bối cảnh thuận lợi cho các sản phẩm, dịch vụ xoay quanh gaming và eSports. Khi cộng đồng người chơi phát triển, nhu cầu về thiết bị gaming như chuột, bàn phím cơ, tai nghe, tay cầm và phụ kiện chuyên dụng cũng tăng theo. Tuy nhiên, không phải người dùng nào cũng có đủ điều kiện hoặc kinh nghiệm để chọn đúng thiết bị phù hợp ngay từ lần mua đầu tiên.

Từ thực tế đó, ý tưởng xây dựng một nền tảng cho phép \textbf{thuê hoặc thử thiết bị gaming trước khi mua} là có tính khả thi. Sản phẩm giải quyết trực tiếp một vấn đề thật của người chơi: gear gaming đắt, khó chọn, và trải nghiệm thực tế quan trọng hơn review trên mạng.

\section{Vấn đề thực tế}

Người chơi game, đặc biệt là người chơi eSports, thường có nhu cầu mua các thiết bị như chuột gaming, bàn phím cơ, tai nghe, tay cầm hoặc các phụ kiện hỗ trợ thi đấu. Tuy nhiên, việc chọn mua gear gaming không đơn giản vì mỗi người có thói quen, kích thước tay, phong cách chơi và cảm nhận khác nhau.

Một số vấn đề thực tế người dùng thường gặp gồm:

\begin{itemize}
    \item \textbf{Gear gaming có giá cao:} nhiều thiết bị chất lượng tốt có giá từ vài trăm nghìn đến vài triệu đồng.
    \item \textbf{Dễ mua sai sản phẩm:} người mới chơi hoặc người chưa có kinh nghiệm rất dễ chọn nhầm chuột, bàn phím, tai nghe hoặc tay cầm không phù hợp.
    \item \textbf{Review không đủ để quyết định:} video review, bài đánh giá và thông số kỹ thuật chỉ cung cấp thông tin tham khảo. Người dùng vẫn cần trải nghiệm thật để biết sản phẩm có phù hợp không.
    \item \textbf{Chuột gaming phụ thuộc vào cảm giác cầm:} một con chuột có thể được đánh giá cao nhưng lại không hợp với kích thước tay, grip style hoặc thói quen DPI của người dùng.
    \item \textbf{Bàn phím cơ phụ thuộc vào loại switch:} người dùng khó biết mình hợp với red switch, brown switch, blue switch, silver switch hay magnetic switch nếu chưa từng thử.
    \item \textbf{Tai nghe gaming cần trải nghiệm lâu:} tai nghe có thể cho âm thanh tốt nhưng đeo lâu bị đau tai, nóng tai hoặc mic không đáp ứng nhu cầu giao tiếp khi chơi game.
    \item \textbf{Tay cầm chơi game cũng cần thử thực tế:} với các game như FC Online, eFootball, FIFA, PES hoặc racing game, độ nhạy nút bấm, độ trễ và cảm giác cầm là yếu tố rất quan trọng.
    \item \textbf{Mua đi bán lại mất thời gian và dễ lỗ:} nếu mua nhầm gear, người dùng thường phải bán lại với giá thấp hơn.
\end{itemize}

Vì vậy, vấn đề cốt lõi là: \textbf{người dùng cần được trải nghiệm thiết bị gaming trước khi bỏ tiền mua chính thức}.

\section{Đối tượng khách hàng}

\subsection{Gamer cá nhân}

Đây là nhóm người chơi game thường xuyên, có nhu cầu nâng cấp thiết bị để cải thiện trải nghiệm chơi game. Họ quan tâm đến hiệu năng, độ thoải mái, độ phản hồi và cảm giác sử dụng thực tế.

\subsection{Người mới chơi eSports}

Nhóm này chưa có nhiều kinh nghiệm chọn gear. Họ dễ bị ảnh hưởng bởi quảng cáo hoặc review trên mạng, từ đó mua sản phẩm không phù hợp với nhu cầu cá nhân.

\subsection{Sinh viên}

Sinh viên là nhóm có nhu cầu chơi game cao nhưng ngân sách hạn chế. Họ có thể không muốn bỏ ra 1--3 triệu đồng để mua một món gear khi chưa chắc sản phẩm đó phù hợp.

\subsection{Người chơi game nghiêm túc hoặc bán chuyên}

Đây là nhóm có nhu cầu tối ưu thiết bị để cải thiện hiệu suất thi đấu. Họ sẵn sàng trả tiền để thử nhiều mẫu gear khác nhau trước khi chọn sản phẩm phù hợp nhất.

\subsection{Esports club, đội tuyển và phòng game}

Các câu lạc bộ hoặc đội tuyển có thể thuê thử nhiều thiết bị cho thành viên trước khi mua số lượng lớn. Phòng game cũng có thể dùng nền tảng để thử nghiệm các dòng gear mới trước khi trang bị cho hệ thống máy.

\subsection{Shop gaming gear}

Các cửa hàng bán thiết bị gaming có thể đưa sản phẩm lên nền tảng để người dùng thuê thử. Đây là một kênh marketing và bán hàng mới, giúp shop tăng độ tiếp cận và tăng khả năng chuyển đổi từ thuê thử sang mua thật.

\section{Giải pháp đưa ra}

Dự án đề xuất xây dựng một nền tảng web/app cho phép người dùng \textbf{thuê thiết bị gaming trong vài ngày hoặc thử trước khi mua}. Người dùng có thể tìm kiếm thiết bị theo loại sản phẩm, thương hiệu, game phù hợp, giá thuê, tình trạng thiết bị và thông số kỹ thuật. Sau đó, người dùng đặt thuê, thanh toán phí thuê, đặt cọc và nhận thiết bị để trải nghiệm thực tế.

\subsection{Các loại thiết bị có thể cho thuê}

\begin{itemize}
    \item Chuột gaming: Logitech, Razer, Zowie, Pulsar, Lamzu, Glorious.
    \item Bàn phím cơ: phân loại theo layout, switch, kết nối và độ phản hồi.
    \item Tai nghe gaming: test âm thanh, mic, độ thoải mái khi đeo lâu.
    \item Tay cầm: phục vụ FC Online, eFootball, FIFA, PES, racing game.
    \item Combo gear: chuột, bàn phím và tai nghe theo từng thể loại game.
\end{itemize}

\subsection{Tính năng cho người thuê}

\begin{itemize}
    \item Đăng ký và đăng nhập tài khoản.
    \item Xem danh sách thiết bị gaming đang có sẵn.
    \item Lọc thiết bị theo loại, thương hiệu, giá thuê, tình trạng và game phù hợp.
    \item Xem chi tiết thiết bị gồm ảnh, thông số, mô tả, tình trạng và đánh giá.
    \item Đặt thuê theo ngày.
    \item Thanh toán phí thuê và đặt cọc giả lập.
    \item Theo dõi trạng thái đơn thuê.
    \item Đánh giá thiết bị sau khi trả hàng.
    \item Nhận gợi ý gear phù hợp theo game đang chơi.
\end{itemize}

\subsection{Tính năng cho shop hoặc chủ gear}

\begin{itemize}
    \item Đăng thiết bị lên nền tảng.
    \item Quản lý danh sách thiết bị.
    \item Cập nhật tình trạng thiết bị: available, booked, renting, returned, maintenance.
    \item Quản lý lịch thuê.
    \item Theo dõi doanh thu.
    \item Xem đánh giá từ người thuê.
\end{itemize}

\subsection{Tính năng cho admin}

\begin{itemize}
    \item Quản lý người dùng.
    \item Quản lý thiết bị.
    \item Quản lý đơn thuê.
    \item Kiểm duyệt thiết bị do shop hoặc chủ gear đăng.
    \item Xử lý tranh chấp, trả muộn, hư hỏng hoặc mất thiết bị.
    \item Theo dõi doanh thu và hoa hồng nền tảng.
\end{itemize}
